
% ============================================================
% 第七章 Gamma 函数及其解析延拓
% 内容：阶乘 → Gamma 函数的积分定义 → 基本性质 → 递推延拓到全平面
% → 余元公式、无零点、斯特林公式（Γ 乘积证明见附录 C）
% 目的：为后续 Zeta 函数提供完整工具
% 编译方式：XeLaTeX（作为 main.tex 的 \input 子文件）
% 注：全书重排后定为第七章（前为第6章「离散部分和与 Perron 积分」；更早曾为第6/5章）
% ============================================================

%==============================================================
\chapter{\texorpdfstring{$\Gamma$}{Gamma} 函数及其解析延拓}\label{ch:gamma_function}
%==============================================================

%==============================================================
\section{引言}\label{sec:gamma_intro}

阶乘 $n!$ 定义在正整数上；$\Gamma$ 函数给出其在复平面上的标准解析延拓，并满足 $\Gamma(n+1)=n!$（$n\in\mathbb{N}$）。
在 $\zeta$ 的积分表示、Hankel 围道与函数方程中，$\Gamma$ 以因子或归一化常数的形式反复出现，因此需要单独成章。

\paragraph{定义域的阶段性}
第~\ref{ch:riemann_zeta_intro}~章在 $\operatorname{Re}s>1$ 上建立 $\zeta$ 的积分表示时，
已临时使用 $\operatorname{Re}s>0$ 上的 Euler 积分 $\Gamma(s)=\int_0^\infty x^{s-1}e^{-x}\,\mathrm{d}x$。
本章对该定义给出收敛性、全纯性与递推的系统讨论，并把定义域从右半平面扩大到
$\mathbb{C}\setminus\{0,-1,-2,\ldots\}$。
若要理解 $\Gamma$ 在 $\operatorname{Re}(s)\le 0$ 的形态（极点、余元、与 $\zeta$ 的函数方程中的因子），
必须完成这一延拓。
解析延拓的定义与唯一性见第~\ref{ch:analytic_continuation_general}~章；本章只处理 $\Gamma$ 本身。

本章安排：先在 $\operatorname{Re}(s)>0$ 上用 Euler 积分定义 $\Gamma$，建立递推；
再由递推完成到 $\mathbb{C}\setminus\{0,-1,-2,\ldots\}$ 的亚纯延拓（极点与留数）；
然后给出余元公式、$\Gamma$ 无零点与 Stirling 渐近，并预览 $\Gamma$ 与 $\zeta$ 的积分联系。
（$1/\Gamma$ 的 Weierstrass 乘积\textbf{完整证明}见附录~\ref{app:weierstrass}，正文主线不依赖该证明。）

\paragraph{历史说明}
Euler（约 1729）给出阶乘的插值表示；Legendre 引入记号 $\Gamma(s)$；Gauss 曾用 $\Pi(s)=\Gamma(s+1)$。本书采用 Legendre 记号。


%==============================================================
\section{从阶乘到积分}\label{sec:factorial_to_integral}
%==============================================================

\subsection{阶乘的积分表示}

首先，我们观察一个重要的积分恒等式。

\begin{theorem}{阶乘的积分表示}{}\label{thm:factorial-integral}
对任意非负整数 $n$，
\begin{equation}
    n! = \int_0^\infty x^n e^{-x} \mathrm{d} x.
    \label{eq:factorial-integral}
\end{equation}
\end{theorem}

\begin{proof}
用数学归纳法证明。当 $n = 0$ 时，约定 $0! = 1$，且
\[
\int_0^\infty e^{-x} \mathrm{d} x = \left[ -e^{-x} \right]_0^\infty = 1.
\]
故 $n = 0$ 时成立。

假设对 $n = k$ 成立，即 $k! = \int_0^\infty x^k e^{-x}\mathrm{d} x$。
考虑 $n = k+1$ 的情况：
利用分部积分，令 $u = x^{k+1}$，$\mathrm{d}v = e^{-x}\mathrm{d} x$，则
$\mathrm{d}u = (k+1)x^k \mathrm{d} x$，$v = -e^{-x}$。于是
\[
\int_0^\infty x^{k+1} e^{-x} \mathrm{d} x
= \left[ -x^{k+1} e^{-x} \right]_0^\infty + (k+1) \int_0^\infty x^k e^{-x} \mathrm{d} x.
\]
边界项：当 $x \to \infty$ 时，$x^{k+1} e^{-x} \to 0$（指数衰减快于任意多项式增长）；
当 $x \to 0^+$ 时，$x^{k+1} e^{-x} \to 0$。因此
\[
\int_0^\infty x^{k+1} e^{-x} \mathrm{d} x = (k+1) \int_0^\infty x^k e^{-x} \mathrm{d} x
= (k+1) \cdot k! = (k+1)!.
\]
由归纳法，公式对所有 $n \ge 0$ 成立。
\end{proof}


这个积分表示是理解 $\Gamma$ 函数的起点：
它将离散的阶乘 $n!$ 表达为一个连续积分。
若定义新函数 $\Gamma(s)$，
使得当 $s = n$（正整数）时 $\Gamma(n) = (n-1)!$，
则可尝试用积分 $\int_0^\infty x^{s-1} e^{-x}\mathrm{d} x$ 来定义。


%==============================================================
\section{\texorpdfstring{$\Gamma$}{Gamma} 函数的积分定义：\texorpdfstring{$\operatorname{Re}(s) > 0$}{Re(s)>0}}\label{sec:gamma_def}
%==============================================================

\subsection{定义与收敛域}

基于上述观察，Euler 将阶乘的积分表示推广到复数。

\begin{definition}{$\Gamma$ 函数，Euler 积分形式}{}\label{def:gamma}
对满足 $\operatorname{Re}(s) > 0$ 的复数 $s$，定义
\begin{equation}
    \boxed{\;
    \Gamma(s) = \int_0^\infty x^{s-1} e^{-x} \mathrm{d} x
    \;\vphantom{\int_0^\infty}.
    }
    \label{eq:gamma-def}
\end{equation}
记号与第~\ref{ch:riemann_zeta_intro}~章积分表示节中的工作定义相同；
该章在 $\operatorname{Re}s>1$ 上推导 $\zeta$ 的积分表示时已临时引用本式，
本章补出收敛性、全纯性、递推与全平面延拓的系统讨论。
\end{definition}


\paragraph{记号说明}
注意定义中的指数是 $x^{s-1}$ 而非 $x^s$。
这样设计的目的是使得 $\Gamma(n) = (n-1)!$，
即正整数处的值比阶乘“滞后”1。
这一约定是历史形成的（Legendre 记号），
虽然有时会带来不便，但已成为标准。


\begin{theorem}{收敛性}{}\label{thm:gamma-convergence}
当 $\operatorname{Re}(s) > 0$ 时，积分 $\int_0^\infty x^{s-1} e^{-x}\mathrm{d} x$ 绝对收敛。
\end{theorem}

\begin{proof}
将积分区间分为两部分：$x \in [0,1]$ 和 $x \in [1, \infty)$。记 $\sigma = \operatorname{Re}(s)$。

\noindent\textbf{在 $x \to 0^+$ 附近：}
对 $x > 0$，有 $x^{s-1} = e^{(s-1)\log x}$，故
\[
|x^{s-1}| = e^{(\sigma-1)\log x} = x^{\sigma-1} \quad (\sigma = \operatorname{Re}(s)),
\]
从而 $|x^{s-1} e^{-x}| = x^{\sigma-1} e^{-x} \sim x^{\sigma-1}$（$x \to 0^+$，因 $e^{-x} \to 1$）。
积分 $\int_0^1 x^{\sigma-1}\mathrm{d} x$ 收敛当且仅当 $\sigma-1 > -1$，
即 $\sigma > 0$。故需 $\operatorname{Re}(s) > 0$。

\noindent\textbf{在 $x \to \infty$ 处：}
指数因子 $e^{-x}$ 起主导作用。对任意实数 $A$，
\[
\lim_{x \to \infty} x^{A} e^{-x} = 0,
\]
因此对任意 $\sigma\in\mathbb{R}$，$\int_1^\infty x^{\sigma-1} e^{-x}\mathrm{d} x$ 绝对收敛。

综上所述，当 $\operatorname{Re}(s) > 0$ 时积分绝对收敛。
\end{proof}

\paragraph{在 $\operatorname{Re}(s)>0$ 上的全纯性}
积分在 $\operatorname{Re}(s)>0$ 上绝对收敛后，由含参量积分的标准结果
（在紧子集上可找可积控制函数，从而可在积分号下对 $s$ 求导，或用 Morera 定理），
$\Gamma(s)$ 在右半平面 $\operatorname{Re}(s)>0$ 上全纯。
含参量积分的全纯性条件见第~\ref{ch:complex_analysis_intro}~章；此处直接使用该结论。

\paragraph{定义域的边界}
在 $\operatorname{Re}(s) = 0$ 时，积分在 $x \to 0^+$ 处发散。
例如 $s = 0$ 时，被积函数 $\sim x^{-1}$，而 $\int_0^1 x^{-1}\mathrm{d} x$ 发散。
因此 $\Gamma(s)$ 的\textbf{积分定义}严格限制在右半平面 $\operatorname{Re}(s) > 0$；
虚轴及左半平面上的值须经下文递推延拓得到。


\subsection{与阶乘的关系}

由定理~\ref{thm:factorial-integral} 直接得到：

\begin{corollary}{}{}\label{cor:gamma-factorial}
对任意正整数 $n$，
\[
\Gamma(n) = (n-1)! = 1 \times 2 \times 3 \times \cdots \times (n-1).
\]
特别地，
\[
\Gamma(1) = 0! = 1, \quad \Gamma(2) = 1! = 1, \quad \Gamma(3) = 2! = 2, \quad \Gamma(4) = 3! = 6.
\]
\end{corollary}

%==============================================================
\section{递推公式}\label{sec:gamma_recurrence}
%==============================================================

\subsection{右半平面上的递推}

$\Gamma$ 函数最重要的性质是满足与阶乘类似的递推关系。

\begin{theorem}{递推公式}{}\label{thm:gamma-recurrence}
对 $\operatorname{Re}(s) > 0$，有
\begin{equation}
    \Gamma(s+1) = s \,\Gamma(s).
    \label{eq:gamma-recurrence}
\end{equation}
\end{theorem}

\begin{proof}
由定义，
\[
\Gamma(s+1) = \int_0^\infty x^{s} e^{-x} \mathrm{d} x.
\]
利用分部积分，令 $u = x^s$，$\mathrm{d}v = e^{-x}\mathrm{d} x$，则
$\mathrm{d}u = s\, x^{s-1}\mathrm{d} x$，$v = -e^{-x}$。于是
\[
\Gamma(s+1) = \left[ -x^s e^{-x} \right]_0^\infty + s \int_0^\infty x^{s-1} e^{-x} \mathrm{d} x.
\]
边界项：当 $x \to \infty$ 时，对任意 $\sigma=\operatorname{Re}(s)$ 有 $x^{\sigma} e^{-x} \to 0$；
当 $x \to 0^+$ 时，$x^s e^{-x} \to 0$（因为 $\operatorname{Re}(s) > 0$）。
因此
\[
\Gamma(s+1) = s \int_0^\infty x^{s-1} e^{-x} \mathrm{d} x = s\,\Gamma(s).
\]
\end{proof}

\subsection{右半平面上的图像}

图~\ref{fig:gamma-real-positive} 给出 $\Gamma(x)$ 在实轴正半轴上的图像（$s=x>0$）。
这是 Euler 积分定义域内的形态，作为下一节向左延拓的对照。

\begin{figure}[htbp]
\centering
\begin{tikzpicture}
\begin{axis}[
    width=12cm,
    height=8cm,
    xlabel={$x$},
    ylabel={$\Gamma(x)$},
    title={延拓前的 $\Gamma$ 函数：仅定义于 $x > 0$},
    xmin=0, xmax=5,
    ymin=0, ymax=12,
    xtick={1,2,3,4,5},
    ytick={2,4,6,8,10,12},
    grid=major,
    grid style={line width=0.3pt, draw=gray!30},
    axis lines=left,
    legend style={at={(0.52,0.72)}, anchor=center, font=\small, fill=white, fill opacity=0.90, draw=gray!40},
    clip=true,
]
% Γ(x)：x→0+ 处上冲示意 →+∞（上端截断于图框）
\addplot[blue!70!black, thick, smooth] coordinates {
    (0.09, 10.616)
    (0.12, 7.7567)
    (0.15, 6.1580)
    (0.18, 5.1660)
    (0.20, 4.5908)
    (0.25, 3.6256)
    (0.30, 2.9916)
    (0.40, 2.2182)
    (0.50, 1.7725)
    (0.60, 1.4892)
    (0.70, 1.2981)
    (0.80, 1.1642)
    (0.90, 1.0686)
    (1.00, 1.0000)
    (1.10, 0.9514)
    (1.20, 0.9182)
    (1.30, 0.8975)
    (1.40, 0.8873)
    (1.46, 0.8856)
    (1.50, 0.8862)
    (1.60, 0.8935)
    (1.70, 0.9086)
    (1.80, 0.9314)
    (1.90, 0.9618)
    (2.00, 1.0000)
    (2.20, 1.1018)
    (2.40, 1.2422)
    (2.60, 1.4296)
    (2.80, 1.6770)
    (3.00, 2.0000)
    (3.20, 2.4240)
    (3.40, 2.9812)
    (3.60, 3.7169)
    (3.80, 4.6942)
    (4.00, 6.0000)
    (4.20, 7.7568)
    (4.50, 11.632)
};
\addplot[red, only marks, mark=*, mark size=2.5pt] coordinates {
    (1,1) (2,1) (3,2) (4,6)
};
\addplot[green!60!black, only marks, mark=triangle*, mark size=3pt] coordinates {
    (1.4616, 0.8856)
};
\legend{$\Gamma(x)$, {$\Gamma(n)=(n-1)!$}, 最小值点}
\end{axis}
\end{tikzpicture}
\caption{延拓前的 $\Gamma(x)$（$x>0$）。
当 $x\to 0^+$ 时 $\Gamma(x)\sim 1/x\to +\infty$；
$x\approx 1.46$ 附近取最小值；整数点 $\Gamma(n)=(n-1)!$。
Euler 积分 $\int_0^\infty e^{-t}t^{x-1}\,\mathrm{d}t$ 仅在 $\operatorname{Re}s>0$ 收敛，
故本图只画实轴正半轴；$\operatorname{Re}s\le 0$ 上的值须经递推延拓给出（见下节与图~\ref{fig:gamma-mod-surface}）。}
\label{fig:gamma-real-positive}
\end{figure}

从图~\ref{fig:gamma-real-positive} 可以看到：
\begin{itemize}
    \item 当 $x \to 0^+$ 时，$\Gamma(x) \to +\infty$（因 $\Gamma(x) \sim 1/x$；图中曲线在左端上冲示意）；
    \item 在 $x \approx 1.4616$ 处取最小值 $\approx 0.8856$（三角点）；
    \item 整数点处 $\Gamma(n) = (n-1)!$（红点；$n=5$ 时 $\Gamma(5)=24$ 已超出图框）；
    \item 当 $x \to \infty$ 时，$\Gamma(x)$ 比任意固定底的指数函数 $a^x$（$a>1$）增长更快
          （Stirling：$\Gamma(x)\sim x^{x-1/2}e^{-x}$）。
\end{itemize}

%==============================================================
\section{解析延拓：从右半平面到全复平面}\label{sec:gamma_continuation}
%==============================================================

上文已在 $\operatorname{Re}(s)>0$ 上定义 $\Gamma$、建立递推，并用图~\ref{fig:gamma-real-positive} 给出正实轴形态。
该半平面之外 Euler 积分不再收敛，故不能沿用同一积分式；
下面用 $\Gamma(s)=\Gamma(s+1)/s$ 将定义推进到左半平面（除极点），
从而完成从右半平面到全平面的延拓。
延拓后在 $s=0,-1,-2,\ldots$ 出现极点，是图~\ref{fig:gamma-real-positive} 无法体现的。

\subsection{用递推公式扩大定义域}

递推公式 $\Gamma(s+1) = s\,\Gamma(s)$ 可以改写为
\begin{equation}
    \Gamma(s) = \frac{\Gamma(s+1)}{s}.
    \label{eq:recurrence-backward}
\end{equation}

这个形式允许我们\textbf{从右半平面向左半平面逐层扩展}。
记第 $k$ 层延拓带为竖条
\[
  -k < \operatorname{Re}(s) \le 1-k
  \qquad (k=1,2,3,\ldots),
\]
并在该带内\textbf{仅排除}本层会碰到的极点 $s=1-k$（即 $s=0,-1,-2,\ldots$ 中落在该带边界上的那一个）。
具体地：

\begin{itemize}
    \item \textbf{积分定义域}：$\operatorname{Re}(s) > 0$，由 Euler 积分给出，全纯。
    \item \textbf{第 1 层}：$-1 < \operatorname{Re}(s) \le 0$，排除 $s=0$；
          此时 $\operatorname{Re}(s+1)>0$，$\Gamma(s+1)$ 已由积分定义，用式~\eqref{eq:recurrence-backward} 定义 $\Gamma(s)$。
    \item \textbf{第 2 层}：$-2 < \operatorname{Re}(s) \le -1$，排除 $s=-1$；
          $\Gamma(s+1)$ 已由第 1 层给出。
    \item \textbf{第 3 层}：$-3 < \operatorname{Re}(s) \le -2$，排除 $s=-2$。
    \item \textbf{第 4 层}：$-4 < \operatorname{Re}(s) \le -3$，排除 $s=-3$。
\end{itemize}
如此继续，可覆盖任意左半竖带。由构造，在相邻层重叠处定义一致，故递推关系 $\Gamma(s+1)=s\Gamma(s)$ 在延拓后的定义域上
（$s,s+1$ 均非极点时）仍然成立。

\begin{figure}[htbp]
    \centering
    \includegraphics[width=0.95\textwidth]{全书图形/6-G模曲面.pdf}
    \caption{$|\Gamma(s)|$ 模曲面：积分定义域 $\operatorname{Re}s>0$ 与四层递推延拓
      （第 $k$ 层对应 $-k<\operatorname{Re}s\le 1-k$，$k=1,2,3,4$；
      金黄色系自右向左逐层加深。竖直虚线与标注标出极点 $s=0,-1,\ldots$）。
      高度为 $|\Gamma|$，在极点处趋于 $+\infty$（图中已裁剪上端）。
      正文可将递推继续到任意左半带；本图只画到第 4 层示意。
      （配图文件名 \texttt{6-G模曲面.pdf} 沿用旧章号编号，内容即本章 $\Gamma$ 模曲面。）}
    \label{fig:gamma-mod-surface}
\end{figure}

\paragraph{图形解读}
图~\ref{fig:gamma-mod-surface} 用\textbf{模} $|\Gamma(s)|$ 展示递推分层，比分别绘制实部、虚部更直接地突出极点与定义域扩展：
\begin{itemize}
    \item 着色自右向左对应：积分域（$\operatorname{Re}s>0$）→ 第 1--4 层延拓带；
    \item 曲面高度为 $|\Gamma(s)|$；在 $s=0,-1,-2,\ldots$ 处急剧升高，对应单极点；
    \item 各竖带在 $\operatorname{Re}s=0,-1,-2,-3$ 等边界处衔接，与 $\Gamma(s)=\Gamma(s+1)/s$ 的构造一致。
\end{itemize}
（严格积分定义域仍是 $\operatorname{Re}s>0$；着色边界为示意。）

如此逐次推进，可将 $\Gamma(s)$ 的定义域扩展到整个复平面，
除 $s = 0, -1, -2, \dots$——在那些点上分母为零而分子 $\Gamma(s+1)$ 有限，故为单极点。

\begin{theorem}{$\Gamma$ 函数的解析延拓}{}\label{thm:gamma-meromorphic}
通过反复应用 $\Gamma(s) = \Gamma(s+1)/s$，可将 $\Gamma$ 从初始定义域 $\operatorname{Re}(s) > 0$
\textbf{延拓}为整个复平面上的亚纯函数：极点恰为 $s = 0, -1, -2, \dots$，均为单极点，留数为
\[
\operatorname{Res}_{s=-n}\, \Gamma(s) = \frac{(-1)^n}{n!}, \qquad n = 0, 1, 2, \dots
\]
（留数证明见下小节。）该延拓在连通区域上的\textbf{唯一性}由解析延拓的唯一性定理保证，
见第~\ref{ch:analytic_continuation_general}~章；此处给出的是显式构造。
\end{theorem}


\paragraph{与解析延拓理论的关系}
上述过程是\textbf{解析延拓}的具体实例：从右半平面上的全纯函数出发，利用函数方程（递推）扩大定义域。
第~\ref{ch:analytic_continuation_general}~章将给出解析延拓的定义与唯一性，并把本节的 $\Gamma$ 递推
抽象为「方法三：函数方程」的一般模式；二者相互印证——本章提供可计算的构造，该章说明其在延拓理论中的位置。


\subsection{极点与留数}

利用递推关系，可以显式写出 $\Gamma$ 在各极点附近的展开。
先看 $s=0$：由 $\Gamma(s)=\Gamma(s+1)/s$，并在 $s=0$ 附近取
$\Gamma(s+1)=\Gamma(1)+\Gamma'(1)s+O(s^2)=1-\gamma s+O(s^2)$
（其中 $\Gamma'(1)=-\gamma$，$\gamma\approx 0.577216$ 为 Euler--Mascheroni 常数），立刻得到\label{prop:gamma-near-zero}
\[
\Gamma(s) = \frac{1}{s} - \gamma + O(s)
\qquad(s\to 0).
\]
对一般负整数 $s=-n$（$n=0,1,2,\dots$），将递推应用 $n+1$ 次：
\[
\Gamma(s) = \frac{\Gamma(s+n+1)}{s(s+1)\cdots(s+n)}.
\]
在 $s=-n$ 附近分母含因子 $(s+n)$，而 $s(s+1)\cdots(s+n-1)$ 非零，
故 $s=-n$ 为单极点。留数为\label{prop:gamma-residues}
\begin{align*}
\operatorname{Res}_{s=-n}\,\Gamma(s)
&= \lim_{s \to -n}\,(s+n)\cdot\frac{\Gamma(s+n+1)}{s(s+1)\cdots(s+n)}
= \frac{\Gamma(1)}{(-n)(-n+1)\cdots(-1)} \\
&= \frac{1}{(-1)^n\,n!}
= \frac{(-1)^n}{n!},
\end{align*}
其中用到 $\Gamma(1)=1$ 与 $(-n)(-n+1)\cdots(-1)=(-1)^n n!$，
以及 $1\big/\!\bigl((-1)^n n!\bigr)=(-1)^n/n!$。

\subsection{递推公式的推广}

由递推公式可以导出更一般的形式：

\begin{corollary}{}{}\label{cor:gamma-shift}
对任意正整数 $n$ 和 $s \notin \{0,-1,-2,\dots,-n+1\}$
（使中间因子均有定义），有
\[
\Gamma(s+n) = s(s+1)(s+2)\cdots(s+n-1) \, \Gamma(s).
\]
特别地，当 $s = 1$ 时，
\[
\Gamma(n+1) = n \cdot (n-1) \cdots 1 \cdot \Gamma(1) = n!.
\]
该式由递推在右半平面上的证明经解析延拓（或直接对延拓后的定义逐次应用）对上述 $s$ 成立。
\end{corollary}

%==============================================================
\section{余元公式（反射公式）}\label{sec:reflection_formula}
%==============================================================

$\Gamma$ 满足将 $\Gamma(s)$ 与 $\Gamma(1-s)$ 联系起来的恒等式（余元公式）。

\begin{theorem}{余元公式}{}\label{thm:reflection}
对 $s \notin \mathbb{Z}$，
\begin{equation}
    \Gamma(s) \,\Gamma(1-s) = \frac{\pi}{\sin \pi s}.
    \label{eq:reflection-formula}
\end{equation}
\end{theorem}

\begin{proof}[证明提要]
\textbf{（1）} 先在带 $0 < \operatorname{Re}(s) < 1$ 上建立等式。
此时 $\Gamma(s)$ 与 $\Gamma(1-s)$ 的 Euler 积分均收敛，且
\[
\bigl|x^{s-1} y^{-s} e^{-(x+y)}\bigr|
= x^{\sigma-1} y^{-\sigma} e^{-(x+y)}
\]
在 $(0,\infty)^2$ 上可积（$\sigma=\operatorname{Re}s\in(0,1)$），故由 Fubini 定理可交换积分次序：
\[
\Gamma(s)\,\Gamma(1-s)
= \int_0^\infty\!\int_0^\infty x^{s-1} y^{-s} e^{-(x+y)}\,\mathrm{d}x\,\mathrm{d}y.
\]
\textbf{（2）} 令 $x=uy$（$u>0$），则 $\mathrm{d}x=y\,\mathrm{d}u$，先对 $y$ 积分得
\[
\Gamma(s)\,\Gamma(1-s)
= \int_0^\infty \frac{u^{s-1}}{1+u}\,\mathrm{d}u.
\]
\textbf{（3）} 经典结果（钥匙孔围道，或沿正实轴割线的 Hankel 型围道，取 $u=-1$ 处留数）给出
\[
\int_0^\infty \frac{u^{s-1}}{1+u}\,\mathrm{d}u = \frac{\pi}{\sin \pi s},
\qquad 0 < \operatorname{Re}(s) < 1.
\]
完整围道计算从略；手法与第~\ref{ch:riemann_zeta_continuation}~章 Hankel 围道相近，亦可查标准复分析教材。
\textbf{（4）} 等式两端作为 $s$ 的亚纯函数在带 $0<\operatorname{Re}s<1$ 上一致，
故由恒等定理（或第~\ref{ch:analytic_continuation_general}~章的解析延拓唯一性）
推广到所有 $s \notin \mathbb{Z}$。
\end{proof}

\begin{corollary}{$\Gamma(1/2)$ 的值}{}\label{cor:gamma-half}
令 $s = 1/2$ 代入余元公式~\eqref{eq:reflection-formula}：
\[
\Gamma\!\left(\tfrac{1}{2}\right)^2 = \frac{\pi}{\sin(\pi/2)} = \pi.
\]
因 $\Gamma$ 在 $(0,\infty)$ 上取正值，取正平方根得
\[
\Gamma\!\left(\frac{1}{2}\right) = \sqrt{\pi}.
\]
\end{corollary}


这是 $\Gamma$ 函数在半整数处取值的来源：
\[
\Gamma\!\left(\frac{3}{2}\right) = \frac{1}{2}\,\Gamma\!\left(\frac{1}{2}\right) = \frac{\sqrt{\pi}}{2},\quad
\Gamma\!\left(\frac{5}{2}\right) = \frac{3}{2} \cdot \frac{1}{2}\,\sqrt{\pi} = \frac{3\sqrt{\pi}}{4},\quad \dots
\]


%==============================================================
\section{\texorpdfstring{$\Gamma$}{Gamma} 无零点；乘积表示（指针）}\label{sec:weierstrass}
%==============================================================

后文函数方程与 $\xi$ 的讨论需要：$\Gamma$ 的零点不会「冒充」$\zeta$ 的零点。
本节用已证的余元公式给出\textbf{无零点}；乘积公式仅作陈述，完整证明放在附录。

\begin{theorem}{$\Gamma$ 无零点}{}\label{cor:gamma-no-zeros}
$\Gamma(s)$ 在 $\mathbb{C}$ 上\textbf{没有零点}：若 $s_0$ 不是 $\Gamma$ 的极点，则 $\Gamma(s_0)\neq 0$
（在极点 $0,-1,-2,\ldots$ 处 $|\Gamma|\to\infty$，亦非有限零点）。
\end{theorem}

\begin{proof}
先设 $s_0\notin\mathbb{Z}$。由余元公式~\eqref{eq:reflection-formula}，
\[
  \Gamma(s_0)\,\Gamma(1-s_0)=\frac{\pi}{\sin(\pi s_0)}.
\]
右端有限且非零。若 $\Gamma(s_0)=0$，则左端为 $0$，矛盾。
故在非整数点 $\Gamma$ 无零点。

再设 $s_0=n$ 为正整数。由 $\Gamma(n)=(n-1)!\neq 0$，无零点。
负整数与 $0$ 已是极点，不是有限零点。证毕。
\end{proof}

\paragraph{乘积表示（非主线必引）}
$\Gamma$ 另有 Weierstrass 形式
\begin{equation}
\frac{1}{\Gamma(s)}
= s\, e^{\gamma s}
\prod_{n=1}^{\infty}\Bigl(1+\frac{s}{n}\Bigr)e^{-s/n}
\label{eq:gamma-weierstrass}
\end{equation}
（$\gamma$ 为 Euler--Mascheroni 常数），以及等价的 Euler 乘积写法。
式~\eqref{eq:gamma-weierstrass} 使 $1/\Gamma$ 的零点（即 $\Gamma$ 的极点）一目了然，
并给出「$1/\Gamma$ 为整函数」的构造性理解。

\textbf{本书正文的证明链不引用}~\eqref{eq:gamma-weierstrass}。
其推导、收敛性及与递推/余元的相容性，作为 Weierstrass 方法的\textbf{完整证明例}，
详见附录~\ref{app:weierstrass}（篇幅不受正文限制，可逐步展开）。
第~\ref{ch:xi_function}~章对 $\xi$ 使用 Hadamard 乘积时，可将本式视为「按极点编码」的对照原型，
而非逻辑前置。

%==============================================================
\section{斯特林公式（渐近展开）}\label{sec:stirling}
%==============================================================

对于大的 $|s|$，$\Gamma$ 函数的行为由斯特林公式描述。

\begin{theorem}{斯特林公式}{}\label{thm:stirling}
当 $|s| \to \infty$ 且 $|\arg s| < \pi - \delta$（$\delta > 0$ 固定），有
\[
\log \Gamma(s) = \left(s - \frac{1}{2}\right)\log s - s + \frac{1}{2}\log(2\pi) + O\!\left(\frac{1}{|s|}\right).
\]
写成指数形式：
\[
\Gamma(s) \sim \sqrt{2\pi} \, s^{s-1/2} e^{-s}.
\]
\end{theorem}

证明从略（可用 Euler--Maclaurin 公式或鞍点法；见标准复分析 / 渐近分析教材）。
本书后文在估计 $\Gamma$ 于竖直线或扇形上的增长时直接引用本式。
特别地，在固定竖直带 $\sigma_1\le\operatorname{Re}s\le\sigma_2$ 上令 $s=\sigma+it$、$|t|\to\infty$，
本式给出 $|\Gamma(\sigma+it)|$ 随 $|t|$ 的衰减（或增长）阶，
供第~\ref{ch:zero_distribution}~、\ref{ch:riemann_explicit_formula}~章中竖线估计使用。

\begin{corollary}{$n!$ 的渐近形式}{}\label{cor:stirling-factorial}
当 $n \to \infty$ 时，
\[
n! \sim \sqrt{2\pi n} \, \left(\frac{n}{e}\right)^n.
\]
\end{corollary}

\begin{proof}
在斯特林公式中令 $s = n$（正整数趋于无穷），得
$\Gamma(n) \sim \sqrt{2\pi}\,n^{n-1/2}e^{-n}$。
由递推公式 $n! = \Gamma(n+1) = n\,\Gamma(n)$，故
\[
n! = n\,\Gamma(n) \sim n \cdot \sqrt{2\pi}\,n^{n-1/2}e^{-n}
= \sqrt{2\pi n}\,n^n e^{-n} = \sqrt{2\pi n}\left(\frac{n}{e}\right)^n. \qedhere
\]
\end{proof}


\paragraph{应用}
斯特林公式用于 $\zeta(s)$ 的函数方程、零点计数 $N(T)$ 以及解析数论中的渐近估计。
具体应用见第~\ref{ch:xi_function}~章（$\xi$ 与 Hadamard 乘积）与第~\ref{ch:zero_distribution}~章（零点几何与 $N(T)$）。


%==============================================================
\section{\texorpdfstring{$\Gamma$}{Gamma} 函数在 \texorpdfstring{$\zeta(s)$}{Zeta(s)} 中的应用预览}\label{sec:gamma_zeta_preview}
%==============================================================

\subsection{积分表示的等价形式}

出发点是 $\Gamma$ 的 Euler 积分定义~\eqref{eq:gamma-def}（$\operatorname{Re}s>0$）：
\[
  \Gamma(s)
  = \int_0^\infty t^{s-1} e^{-t}\,\mathrm{d}t.
\]
对此积分作代换 $t=nx$，其中参数 $n$ 满足 $\operatorname{Re}n>0$（实轴情形取 $n>0$ 即可），
则 $\mathrm{d}t=n\,\mathrm{d}x$，积分限仍为 $(0,\infty)$，从而得到缩放形式
\begin{equation}
\int_0^\infty x^{s-1} e^{-nx} \,\mathrm{d}x = \frac{\Gamma(s)}{n^s},
\qquad \operatorname{Re}(s) > 0,\; \operatorname{Re}n > 0.
\label{eq:gamma-scaled}
\end{equation}
（$n$ \textbf{不必}为正整数；对 $\zeta$ 求和时再限制 $n=1,2,3,\ldots$。）

\begin{proof}
在~\eqref{eq:gamma-def} 中令 $t=nx$，则 $x=t/n$，$\mathrm{d}t=n\,\mathrm{d}x$
（$\operatorname{Re}n>0$ 时路径仍为 $(0,\infty)$）。为与~\eqref{eq:gamma-scaled} 左端一致，
将哑变量仍记作 $x$，即
\[
\Gamma(s)
= \int_0^\infty (nx)^{s-1} e^{-nx}\, n\,\mathrm{d}x
= n^{s}\int_0^\infty x^{s-1} e^{-nx}\,\mathrm{d}x,
\]
两边除以 $n^{s}$ 即得~\eqref{eq:gamma-scaled}。\qedhere
\end{proof}

\subsection{\texorpdfstring{$\zeta(s)$}{Zeta(s)} 的 \texorpdfstring{$\Gamma$}{Gamma} 积分表示（复习）}

第~\ref{ch:riemann_zeta_intro}~章已给出：对 $\operatorname{Re}(s) > 1$，
\begin{equation}
    \boxed{\;
    \zeta(s) = \frac{1}{\Gamma(s)} \int_0^\infty \frac{x^{s-1}}{e^x - 1} \,\mathrm{d}x
    \;\vphantom{\int_0^\infty}.
    }
    \label{eq:zeta-gamma-integral}
\end{equation}
推导要点：对每个\textbf{正整数} $n\ge 1$ 应用~\eqref{eq:gamma-scaled} 得
$\Gamma(s)/n^s=\int_0^\infty x^{s-1}e^{-nx}\,\mathrm{d}x$，
再对 $\zeta(s)=\sum_{n=1}^\infty n^{-s}$ 求和，并在 $\operatorname{Re}s>1$ 下由\textbf{绝对收敛}交换求和与积分，得
$\Gamma(s)\zeta(s)=\int_0^\infty x^{s-1}/(e^x-1)\,\mathrm{d}x$。

此处 $\Gamma(s)$ 起归一化因子的作用；式~\eqref{eq:zeta-gamma-integral} 仍限于 $\operatorname{Re}s>1$，
将作为第~\ref{ch:riemann_zeta_continuation}~章延拓 $\zeta$ 与推导函数方程的出发点。

%==============================================================
\section*{延伸阅读}
\addcontentsline{toc}{section}{延伸阅读}
%==============================================================

以下内容为 $\Gamma$ 函数的补充知识，不影响后续章节的理解，
感兴趣的读者可深入阅读。

\subsection*{$\Gamma$ 函数与 Beta 函数的关系}

Beta 函数定义为
\[
\mathrm{B}(p, q) = \int_0^1 t^{p-1} (1-t)^{q-1}\mathrm{d} t, \qquad \operatorname{Re}(p)>0,\; \operatorname{Re}(q)>0.
\]
它与 $\Gamma$ 函数满足恒等式
\[
\mathrm{B}(p, q) = \frac{\Gamma(p)\,\Gamma(q)}{\Gamma(p+q)}.
\]

这一关系可用于推导余元公式、计算许多涉及 Beta 函数的积分，
也给出 $\Gamma$ 函数与组合数 $\binom{n}{k}$ 推广形式之间的联系。
详细证明见标准复分析或特殊函数教材（本书附录不收录 Beta 函数专论）。

\begin{center}
\fbox{\parbox{0.92\textwidth}{
\centering
\vspace{0.2em}
\textbf{第七章：重要定理与核心公式汇总}\\[6pt]
\renewcommand{\arraystretch}{1.6}
\begin{tabular}{lp{9.5cm}}
\hline
\textbf{名称} & \textbf{数学内容 / 公式表达} \\
\hline
$\Gamma$ 函数定义 & $\Gamma(s)=\displaystyle\int_0^{\infty} t^{s-1}e^{-t}\,\mathrm{d}t$，$\Re(s)>0$（全纯） \\
递推与延拓 & $\Gamma(s+1)=s\,\Gamma(s)$；递推向左延拓为 $\mathbb{C}$ 上亚纯函数 \\
极点与无零点 & 极点 $0,-1,-2,\ldots$（单极点）；$\Gamma(s)\neq 0$（余元，定理~\ref{cor:gamma-no-zeros}） \\
余元公式 & $\Gamma(s)\Gamma(1-s)=\dfrac{\pi}{\sin\pi s}$（$s\notin\mathbb{Z}$） \\
Weierstrass 乘积（附录） &
  $\dfrac{1}{\Gamma(s)}=s\,e^{\gamma s}\displaystyle\prod_{n=1}^{\infty}\bigl(1+\dfrac{s}{n}\bigr)e^{-s/n}$；
  证明见附录~\ref{app:weierstrass}，正文主线不引 \\
Stirling 公式 & $\Gamma(s)\sim \sqrt{2\pi}\,s^{s-1/2}e^{-s}$，
  $|s|\to\infty$ 且 $|\arg s|<\pi-\delta$ \\
\hline
\end{tabular}
\vspace{0.2em}
}}
\end{center}

\section*{本章小结与后续展望}\label{sec:gamma_summary}
\addcontentsline{toc}{section}{本章小结与后续展望}

本章在 $\operatorname{Re}s>0$ 上用 Euler 积分定义 $\Gamma$
（与第~\ref{ch:riemann_zeta_intro}~章工作定义一致），证明收敛与全纯及递推，
再由 $\Gamma(s)=\Gamma(s+1)/s$ \textbf{构造}到全平面的亚纯延拓
（构造在本章完成；延拓在连通区域上的\textbf{唯一性}见第~\ref{ch:analytic_continuation_general}~章，
递推本身并不代替唯一性定理）。
极点恰为 $0,-1,-2,\ldots$，均为单极点；由余元知 $\Gamma$ 全平面无零点。
余元与 Stirling 给出进一步结构；$1/\Gamma$ 的 Weierstrass 乘积证明见附录~\ref{app:weierstrass}（选读）。

$\Gamma$ 出现在 $\zeta$ 的积分表示与函数方程中。下一章给出解析延拓的定义与唯一性；
第~\ref{ch:riemann_zeta_continuation}~章再具体处理 $\zeta$ 的延拓与函数方程。